\documentclass[10pt,a4paper]{amsart}
\usepackage{amsmath,amssymb,amsthm,graphicx}
\usepackage{verbatim}
\newenvironment{reference}{\comment}{\endcomment}
%\newenvironment{reference}{}{}
\newenvironment{slogan}{\comment}{\endcomment}
\newenvironment{history}{\comment}{\endcomment}

% For commutative diagrams we use Xy-pic
\usepackage[all]{xy}

% We use 2cell for 2-commutative diagrams.
\xyoption{2cell}
% Theorem environments.
%
\theoremstyle{plain}
\newtheorem{theorem}[subsection]{Theorem}
\newtheorem{proposition}[subsection]{Proposition}
\newtheorem{lemma}[subsection]{Lemma}

\theoremstyle{definition}
\newtheorem{definition}[subsection]{Definition}
\newtheorem{example}[subsection]{Example}
\newtheorem{exercise}[subsection]{Exercise}
\newtheorem{situation}[subsection]{Situation}

\theoremstyle{remark}
\newtheorem{remark}[subsection]{Remark}
\newtheorem{remarks}[subsection]{Remarks}

\numberwithin{equation}{subsection}

% Macros
%
\def\lim{\mathop{\mathrm{lim}}\nolimits}
\def\colim{\mathop{\mathrm{colim}}\nolimits}
\def\Spec{\mathop{\mathrm{Spec}}}
\def\Hom{\mathop{\mathrm{Hom}}\nolimits}
\def\Ext{\mathop{\mathrm{Ext}}\nolimits}
\def\SheafHom{\mathop{\mathcal{H}\!\mathit{om}}\nolimits}
\def\SheafExt{\mathop{\mathcal{E}\!\mathit{xt}}\nolimits}
\def\Sch{\mathit{Sch}}
\def\Mor{\mathop{\mathrm{Mor}}\nolimits}
\def\Ob{\mathop{\mathrm{Ob}}\nolimits}
\def\Sh{\mathop{\mathit{Sh}}\nolimits}
\def\NL{\mathop{N\!L}\nolimits}
\def\CH{\mathop{\mathrm{CH}}\nolimits}
\def\proetale{{pro\text{-}\acute{e}tale}}
\def\etale{{\acute{e}tale}}
\def\QCoh{\mathit{QCoh}}
\def\Ker{\mathop{\mathrm{Ker}}}
\def\Im{\mathop{\mathrm{Im}}}
\def\Coker{\mathop{\mathrm{Coker}}}
\def\Coim{\mathop{\mathrm{Coim}}}

% Boxtimes
%
\DeclareMathSymbol{\boxtimes}{\mathbin}{AMSa}{"02}

%
% Macros for moduli stacks/spaces
%
\def\QCohstack{\mathcal{QC}\!\mathit{oh}}
\def\Cohstack{\mathcal{C}\!\mathit{oh}}
\def\Spacesstack{\mathcal{S}\!\mathit{paces}}
\def\Quotfunctor{\mathrm{Quot}}
\def\Hilbfunctor{\mathrm{Hilb}}
\def\Curvesstack{\mathcal{C}\!\mathit{urves}}
\def\Polarizedstack{\mathcal{P}\!\mathit{olarized}}
\def\Complexesstack{\mathcal{C}\!\mathit{omplexes}}
% \Pic is the operator that assigns to X its picard group, usage \Pic(X)
% \Picardstack_{X/B} denotes the Picard stack of X over B
% \Picardfunctor_{X/B} denotes the Picard functor of X over B
\def\Pic{\mathop{\mathrm{Pic}}\nolimits}
\def\Picardstack{\mathcal{P}\!\mathit{ic}}
\def\Picardfunctor{\mathrm{Pic}}
\def\Deformationcategory{\mathcal{D}\!\mathit{ef}}
\usepackage{cleveref}
\crefname{originallemma}{original lemma}{original lemmas}
\Crefname{originallemma}{Original lemma}{Original lemmas}
\makeatletter
\newlabel{morphisms-lemma-generalizations-lift-flat}{{03HV}{source}{}{sourcepointer-1}{}}
\newlabel{morphisms-lemma-generalizations-lift-flat@cref}{{[originallemma][0][]03HV}{[1][0][]source}{}{}{}}
\makeatother
\crefname{originalreference}{original reference}{original references}
\Crefname{originalreference}{Original reference}{Original references}
\makeatletter
\newlabel{situation-pre-pure}{{05FW}{source}{}{sourcepointer-2}{}}
\newlabel{situation-pre-pure@cref}{{[originalreference][0][]05FW}{[1][0][]source}{}{}{}}
\makeatother
\crefname{originallemma}{original lemma}{original lemmas}
\Crefname{originallemma}{Original lemma}{Original lemmas}
\makeatletter
\newlabel{spaces-divisors-lemma-base-change-relative-assassin}{{0CV2}{source}{}{sourcepointer-3}{}}
\newlabel{spaces-divisors-lemma-base-change-relative-assassin@cref}{{[originallemma][0][]0CV2}{[1][0][]source}{}{}{}}
\makeatother
\crefname{originallemma}{original lemma}{original lemmas}
\Crefname{originallemma}{Original lemma}{Original lemmas}
\makeatletter
\newlabel{spaces-morphisms-lemma-base-change-flat}{{03MO}{source}{}{sourcepointer-4}{}}
\newlabel{spaces-morphisms-lemma-base-change-flat@cref}{{[originallemma][0][]03MO}{[1][0][]source}{}{}{}}
\makeatother
\crefname{originallemma}{original lemma}{original lemmas}
\Crefname{originallemma}{Original lemma}{Original lemmas}
\makeatletter
\newlabel{spaces-properties-lemma-points-cartesian}{{03H4}{source}{}{sourcepointer-5}{}}
\newlabel{spaces-properties-lemma-points-cartesian@cref}{{[originallemma][0][]03H4}{[1][0][]source}{}{}{}}
\makeatother
\title{Impurities ascend along a flat base change}
\author{The Stacks Project Authors}
\date{stacks-project:spaces-flat.tex:local-snapshot-20261012; GNU FDL 1.2 or later}
\begin{document}
\maketitle
\noindent\textit{Editorial note.} All delivered mathematical text is the author's own, copied from the pinned Stacks Project chapter file spaces-flat.tex; no statement or proof step was rewritten, reordered or removed.\par\medskip
\noindent\textit{Editorial note.} The counted proof is the author's complete argument: an etale chart construction of the fibre product, a lift of the specialisation along the flat induced morphism, and a closure comparison using base change of the relative assassin\par\medskip
\noindent\textit{Editorial note.} The label lemma-flat-ascent-impurity names the counted statement, and the proof environment immediately following it in the source is delivered verbatim as its proof.\par\medskip
\noindent\textit{Editorial note.} The counted statement is carried by the author's own lemma environment in the chapter source.\par\medskip
\noindent\textit{Allowed standard prerequisites (editorial).}\begin{itemize}
\item Impurities above a point, and the base change of the relative assassin of a module
\item The existence of an etale chart of a fibre product through a prescribed pair of points, and the resulting flatness of the induced morphism of schemes
\item Lifting of specialisations along a flat morphism of schemes, and the compatibility of the closed subspace cut out by a point with the base change
\end{itemize}\medskip
\ifdefined\hypertarget\hypertarget{sourcepointer-1}{}\fi
\noindent\textit{Source reference.} Original lemma 03HV: sources/stacks-chow/tags.txt, source lines 4581--4581. Prerequisite pointer invoked at proof line(s) 387; it names a standard fact of a sibling chapter that the counted argument uses, and every step of the counted proof built on it is retained verbatim in the delivered body.\par\medskip
\ifdefined\hypertarget\hypertarget{sourcepointer-2}{}\fi
\noindent\textit{Source reference.} Original reference 05FW: sources/stacks-spaces-flat/pinned-original-source.tex, source lines 260--260. Prerequisite pointer invoked at proof line(s) 367; it names a standard fact of the same chapter that the counted argument uses, and every step of the counted proof built on it is retained verbatim in the delivered body.\par\medskip
\ifdefined\hypertarget\hypertarget{sourcepointer-3}{}\fi
\noindent\textit{Source reference.} Original lemma 0CV2: sources/stacks-chow/tags.txt, source lines 16122--16122. Prerequisite pointer invoked at proof line(s) 393; it names a standard fact of a sibling chapter that the counted argument uses, and every step of the counted proof built on it is retained verbatim in the delivered body.\par\medskip
\ifdefined\hypertarget\hypertarget{sourcepointer-4}{}\fi
\noindent\textit{Source reference.} Original lemma 03MO: sources/stacks-chow/tags.txt, source lines 4756--4756. Prerequisite pointer invoked at proof line(s) 383; it names a standard fact of a sibling chapter that the counted argument uses, and every step of the counted proof built on it is retained verbatim in the delivered body.\par\medskip
\ifdefined\hypertarget\hypertarget{sourcepointer-5}{}\fi
\noindent\textit{Source reference.} Original lemma 03H4: sources/stacks-chow/tags.txt, source lines 4553--4553. Prerequisite pointer invoked at proof line(s) 381; it names a standard fact of a sibling chapter that the counted argument uses, and every step of the counted proof built on it is retained verbatim in the delivered body.\par\medskip
\begin{lemma}
\label{lemma-flat-ascent-impurity}
In Situation \ref{situation-pre-pure}.
Let $(Y_1, y_1) \to (Y, y)$ be a morphism of pointed
algebraic spaces over $S$. Assume $Y_1 \to Y$ is flat at $y_1$.
If $(T \to Y, t' \leadsto t, \xi)$ is an impurity of
$\mathcal{F}$ above $y$, then there exists an impurity
$(T_1 \to Y_1, t_1' \leadsto t_1, \xi_1)$ of the pullback
$\mathcal{F}_1$ of $\mathcal{F}$ to $X_1 = Y_1 \times_Y X$
over $y_1$ such that $T_1$ is \'etale over $Y_1 \times_Y T$.
\end{lemma}

\begin{proof}
Choose an \'etale morphism $T_1 \to Y_1 \times_Y T$ where $T_1$
is a scheme and let $t_1 \in T_1$ be a point mapping to $y_1$ and $t$.
It is possible to find a pair $(T_1, t_1)$ like this by
Properties of Spaces, Lemma \ref{spaces-properties-lemma-points-cartesian}.
The morphism of schemes $T_1 \to T$ is flat at $t_1$
(use Morphisms of Spaces, Lemma \ref{spaces-morphisms-lemma-base-change-flat}
and the definition of flat morphisms of algebraic spaces)
there exists a specialization $t'_1 \leadsto t_1$ lying over
$t' \leadsto t$, see
Morphisms, Lemma \ref{morphisms-lemma-generalizations-lift-flat}.
Choose a point $\xi_1 \in |X_{T_1}|$ mapping to $t'_1$
and $\xi$ with $\xi_1 \in \text{Ass}_{X_{T_1}/T_1}(\mathcal{F}_{T_1})$.
point of $\Spec(\kappa(t'_1) \otimes_{\kappa(t')} \kappa(\xi))$.
This is possible by
Divisors on Spaces, Lemma
\ref{spaces-divisors-lemma-base-change-relative-assassin}.
As the closure $Z_1$ of $\{\xi_1\}$ in $|X_{T_1}|$ maps into the
closure of $\{\xi\}$ in $|X_T|$ we conclude that
the image of $Z_1$ in $|T_1|$ cannot contain $t_1$.
Hence $(T_1 \to Y_1, t'_1 \leadsto t_1, \xi_1)$
is an impurity of $\mathcal{F}_1$ above $Y_1$.
\end{proof}

\end{document}
